%% schriftsatz.latex — Anpassbares Pandoc-Template für familienrechtliche Schriftsätze
%%
%% Pandoc-Variablen (via -V oder YAML-Frontmatter):
%%   mainfont      Schriftart             (Standard: Arial)
%%   fontsize      Schriftgröße           (Standard: 11pt)
%%   papersize     Papierformat           (Standard: a4)
%%   geometry      Seitenränder           (Standard: margin=2.5cm)
%%   linestretch   Zeilenabstand          (Standard: 1.2)
%%   lang          Sprache                (Standard: de-DE)
%%   briefkopf     Vorbereiteter LaTeX-Briefkopf (von generate-pdf.py erzeugt)
%%
%% ─────────────────────────────────────────────────────────────────────────────

\documentclass[
  11pt,
  a4paper,
]{article}

%% ── Sprache ──────────────────────────────────────────────────────────────────
\usepackage[ngerman]{babel}

%% ── Fonts (XeLaTeX) ──────────────────────────────────────────────────────────
\usepackage{fontspec}
\setmainfont{Arial}[
  BoldFont       = Arial Bold,
  ItalicFont     = Arial Italic,
  BoldItalicFont = Arial Bold Italic,
  Ligatures      = TeX,
]
\usepackage{microtype}

%% ── Seitenlayout ─────────────────────────────────────────────────────────────
\usepackage[left=3cm,right=2cm,top=2.5cm,bottom=2.5cm]{geometry}

%% ── Zeilenabstand & Absätze ──────────────────────────────────────────────────
\usepackage{setspace}
\setstretch{1.3}
\usepackage{parskip}
\setlength{\parskip}{6pt}
\setlength{\parindent}{0pt}

%% ── Überschriften ────────────────────────────────────────────────────────────
\usepackage{titlesec}
\titleformat{\section}{\normalfont\large\bfseries}{}{0em}{}
\titleformat{\subsection}{\normalfont\normalsize\bfseries}{}{0em}{}
\titleformat{\subsubsection}{\normalfont\normalsize\bfseries}{}{0em}{}
\titlespacing{\section}{0pt}{14pt}{6pt}
\titlespacing{\subsection}{0pt}{10pt}{4pt}
%% Keine automatische Abschnittsnummerierung
\setcounter{secnumdepth}{0}

%% ── Seitenzahlen ─────────────────────────────────────────────────────────────
\usepackage{fancyhdr}
\usepackage{lastpage}
\pagestyle{fancy}
\fancyhf{}
\fancyfoot[C]{\footnotesize Seite~\thepage~von~\pageref{LastPage}}
\renewcommand{\headrulewidth}{0pt}
\renewcommand{\footrulewidth}{0pt}

%% ── Tabellen ─────────────────────────────────────────────────────────────────
\usepackage{booktabs}
\usepackage{longtable}
\usepackage{array}
\usepackage{calc}
\usepackage{caption}
\captionsetup{font=small, labelfont=bf}
%% Pandoc 3.x setzt \LTcaptype auf "none" — LaTeX braucht diesen Counter:
\newcounter{none}

%% ── Listen ───────────────────────────────────────────────────────────────────
\usepackage{enumitem}
\setlist{noitemsep, topsep=4pt}

%% ── Sonstiges ────────────────────────────────────────────────────────────────
\usepackage{xcolor}
\usepackage[normalem]{ulem}
\usepackage[hidelinks]{hyperref}
\usepackage{needspace}

%% ── Pandoc-Kompatibilität ────────────────────────────────────────────────────
\providecommand{\tightlist}{%
  \setlength{\itemsep}{0pt}\setlength{\parskip}{0pt}}

%% ─────────────────────────────────────────────────────────────────────────────
\begin{document}

%% Briefkopf — von generate-pdf.py aus Markdown-Zeilen vor dem ersten --- erzeugt.
%% Zeilen mit "den \d" werden rechtsbündig gesetzt, Betreff fett+unterstrichen,
%% Aktenzeichen fett, alle anderen Zeilen linksbündig mit explizitem Zeilenumbruch.

\section{Fakten --- Berger ./. Berger}\label{fakten-berger-..-berger}

\subsection{Verfahrensdaten}\label{verfahrensdaten}

{\def\LTcaptype{none} % do not increment counter
\begin{longtable}[]{@{}ll@{}}
\toprule\noalign{}
Feld & Wert \\
\midrule\noalign{}
\endhead
\bottomrule\noalign{}
\endlastfoot
Aktenzeichen & 7 F 88/25 \\
Gericht & Amtsgericht Neukirchen \\
Antragstellerin & Sandra Berger \\
Anwalt/Anwältin AS & RA Petra Schmitt \\
Antragsgegner & Thomas Berger \\
Kind(er) & Finn Berger, geb. 03.05.2019 \\
Verfahrensbeistand & Dipl.-Soz. Klaus Hofmann \\
Beratungsstelle & AWO Beratungsstelle Neukirchen \\
\end{longtable}
}

\subsection{Persönliche Daten
Mandant}\label{persuxf6nliche-daten-mandant}

{\def\LTcaptype{none} % do not increment counter
\begin{longtable}[]{@{}
  >{\raggedright\arraybackslash}p{(\linewidth - 2\tabcolsep) * \real{0.5000}}
  >{\raggedright\arraybackslash}p{(\linewidth - 2\tabcolsep) * \real{0.5000}}@{}}
\toprule\noalign{}
\begin{minipage}[b]{\linewidth}\raggedright
Feld
\end{minipage} & \begin{minipage}[b]{\linewidth}\raggedright
Wert
\end{minipage} \\
\midrule\noalign{}
\endhead
\bottomrule\noalign{}
\endlastfoot
Name & Thomas Berger \\
Beruf & Schreinermeister \\
Arbeitszeitmodell & Vollzeit, Mo--Fr \\
Wohnsituation & 3-Zimmer-Wohnung, Tannenweg 18, 87654 Neukirchen, 76
qm \\
Entfernung zur Wohnung des anderen ET & 1,8 km (ca. 6 Min. mit dem
Rad) \\
\end{longtable}
}

\subsection{Betreuungsrealität}\label{betreuungsrealituxe4t}

\subsubsection{Vor der Trennung}\label{vor-der-trennung}

\begin{itemize}
\tightlist
\item
  Morgenversorgung: überwiegend gemeinsam; bei Frühschicht AG allein
  (ca. 3×/Woche)
\item
  KiTa-Bringen/-Abholen: Bringen ca. 60\% AG, Abholen ca. 50/50
\item
  Zubettgehen: abwechselnd, je nach Arbeitsende; AG übernahm ca.
  4×/Woche
\item
  Freizeitaktivitäten: AG organisierte Schwimmkurs und Bastelgruppe
\item
  Haushaltshilfe/Tagesmutter: keine; KiTa ganztags ab 07:30 Uhr
\end{itemize}

\subsubsection{Nach der Trennung}\label{nach-der-trennung}

\begin{itemize}
\tightlist
\item
  Aktuelles Betreuungsmodell: {[}x{]} Paritätisches Wechselmodell
\item
  Aufteilung in Prozent (Antragsgegner / Antragstellerin): 50 / 50
\item
  Seit wann: 01.09.2025
\item
  Vereinbart oder faktisch gelebt: schriftlich vereinbart
  (Elternvereinbarung 01.09.2025, Anlage B1)
\item
  Arbeitszeitreduzierung: keine
\item
  Unterstützung durch Arbeitgeber: flexible Arbeitszeiten in der eigenen
  Schreinerei
\item
  Angestrebtes Modell (eigene Position): Beibehaltung paritätisches
  Wechselmodell
\end{itemize}

\subsection{Kernargumente (eigene
Stärken)}\label{kernargumente-eigene-stuxe4rken}

\begin{enumerate}
\def\labelenumi{\arabic{enumi}.}
\tightlist
\item
  Aktive Beteiligung seit Geburt belegt durch KiTa-Anwesenheitslisten
  und Betreuungsprotokoll
\item
  Stabile Wohnsituation: kindgerechte 3-Zimmer-Wohnung mit eigenem
  Zimmer für Finn, gleicher KiTa-Bezirk
\item
  KiTa-Bestätigung (Entwicklungsgespräch 25.02.2026): Finn entwickelt
  sich altersgemäß, keine Verhaltensauffälligkeiten
\end{enumerate}

\subsection{Behauptungen der
Gegenseite}\label{behauptungen-der-gegenseite}

{\def\LTcaptype{none} % do not increment counter
\begin{longtable}[]{@{}
  >{\raggedright\arraybackslash}p{(\linewidth - 10\tabcolsep) * \real{0.0909}}
  >{\raggedright\arraybackslash}p{(\linewidth - 10\tabcolsep) * \real{0.2000}}
  >{\raggedright\arraybackslash}p{(\linewidth - 10\tabcolsep) * \real{0.1636}}
  >{\raggedright\arraybackslash}p{(\linewidth - 10\tabcolsep) * \real{0.1818}}
  >{\raggedright\arraybackslash}p{(\linewidth - 10\tabcolsep) * \real{0.2000}}
  >{\raggedright\arraybackslash}p{(\linewidth - 10\tabcolsep) * \real{0.1636}}@{}}
\toprule\noalign{}
\begin{minipage}[b]{\linewidth}\raggedright
Nr.
\end{minipage} & \begin{minipage}[b]{\linewidth}\raggedright
Behauptung
\end{minipage} & \begin{minipage}[b]{\linewidth}\raggedright
Schwere
\end{minipage} & \begin{minipage}[b]{\linewidth}\raggedright
Belegbar
\end{minipage} & \begin{minipage}[b]{\linewidth}\raggedright
Strategie
\end{minipage} & \begin{minipage}[b]{\linewidth}\raggedright
Antwort
\end{minipage} \\
\midrule\noalign{}
\endhead
\bottomrule\noalign{}
\endlastfoot
1 & AG war berufsbedingt häufig abwesend und konnte keine verlässliche
Betreuung sicherstellen & H & J & Widerlegen & Dienstreisen betrugen im
gesamten Jahr 2025 max. 10 Tage; Hotelrechnungen und Kalender belegen
genaue Daten (B4) \\
2 & Elternvereinbarung wurde unter emotionalem Druck unterzeichnet & M &
N & Einordnen & Vereinbarung wurde über mehrere Wochen gemeinsam
ausgehandelt; AWO-Beratung bestätigt freiwillige Einigung \\
\end{longtable}
}

\subsection{Widersprüche im Antrag der
Gegenseite}\label{widerspruxfcche-im-antrag-der-gegenseite}

{\def\LTcaptype{none} % do not increment counter
\begin{longtable}[]{@{}
  >{\raggedright\arraybackslash}p{(\linewidth - 6\tabcolsep) * \real{0.1250}}
  >{\raggedright\arraybackslash}p{(\linewidth - 6\tabcolsep) * \real{0.3000}}
  >{\raggedright\arraybackslash}p{(\linewidth - 6\tabcolsep) * \real{0.3000}}
  >{\raggedright\arraybackslash}p{(\linewidth - 6\tabcolsep) * \real{0.2750}}@{}}
\toprule\noalign{}
\begin{minipage}[b]{\linewidth}\raggedright
Nr.
\end{minipage} & \begin{minipage}[b]{\linewidth}\raggedright
Widerspruch
\end{minipage} & \begin{minipage}[b]{\linewidth}\raggedright
Fundstelle
\end{minipage} & \begin{minipage}[b]{\linewidth}\raggedright
Bedeutung
\end{minipage} \\
\midrule\noalign{}
\endhead
\bottomrule\noalign{}
\endlastfoot
1 & Antrag behauptet Reisetag AG am 15.10.2025, AG war nachweislich in
Neukirchen & P7 vs B4 & Belegt selektive rückwirkende
Protokollerstellung durch AS \\
\end{longtable}
}

\subsection{Belege-Inventar}\label{belege-inventar}

{\def\LTcaptype{none} % do not increment counter
\begin{longtable}[]{@{}
  >{\raggedright\arraybackslash}p{(\linewidth - 6\tabcolsep) * \real{0.2222}}
  >{\raggedright\arraybackslash}p{(\linewidth - 6\tabcolsep) * \real{0.2222}}
  >{\raggedright\arraybackslash}p{(\linewidth - 6\tabcolsep) * \real{0.2222}}
  >{\raggedright\arraybackslash}p{(\linewidth - 6\tabcolsep) * \real{0.3333}}@{}}
\toprule\noalign{}
\begin{minipage}[b]{\linewidth}\raggedright
Anlage
\end{minipage} & \begin{minipage}[b]{\linewidth}\raggedright
Inhalt
\end{minipage} & \begin{minipage}[b]{\linewidth}\raggedright
Status
\end{minipage} & \begin{minipage}[b]{\linewidth}\raggedright
Beweiswert
\end{minipage} \\
\midrule\noalign{}
\endhead
\bottomrule\noalign{}
\endlastfoot
B1 & Elternvereinbarung vom 01.09.2025 (paritätisches Wechselmodell) &
vorhanden & hoch \\
B2 & KiTa-Bericht Regenbogen, Halbjahresbericht Herbst 2025 & vorhanden
& hoch \\
B3 & WhatsApp-Verlauf AS/AG Oktober--November 2025
(Übergabe-Koordination) & vorhanden & mittel \\
B4 & Dienstreise-Belege: Hotelrechnung Stuttgart 16.10.2025, Zugtickets
Düsseldorf 18.--20.02.2026 & vorhanden & hoch \\
B5 & E-Mail KiTa-Entwicklungsgespräch 10.02.2026 (AS in CC) & vorhanden
& hoch \\
B6 & Gemeinsamer Jahreskalender 2025/2026 (Google-Calendar-Export) &
vorhanden & mittel \\
\end{longtable}
}

\end{document}
